% Complete Section 7.8 and Box 7.1, original PDF 103-106 / printed 93-96.
\endgroup
\begingroup\emergencystretch=2em
% block: historical-antecedents-heading page:103 type:section
\section{{历史渊源}}\label{sec:historical-antecedents}

% block: finite-element-families-unification page:103 type:ordinary-paragraph
{对任意维数 $n$, 每个 $0\leq k\leq n$ 和 $r\geq1$, 已用系统而统一的方式定义了单纯形有限元空间 $\mathcal P_r\Lambda^k(\mathcal T_h)$ 和 $\mathcal P_r^-\Lambda^k(\mathcal T_h)$. 在一维, 二维和三维中, 这些有限元微分形式空间可通过代理与标量值和向量值有限元空间等同. 这样, 它们统一了许多研究者在数十年间相对独立地发展出的一大类有限元方法.}

% block: classical-finite-element-antecedents pages:103,104 type:ordinary-paragraph
{在任意维数中, 空间 $\mathcal P_r\Lambda^0$ 和 $\mathcal P_r^-\Lambda^0$ 彼此重合, 并与最经典的有限元空间等同, 即次数至多为 $r$ 的连续分片多项式构成的 Lagrange 空间; 当 $r=1$ 时, 它可追溯到 Courant 于 1943 年的工作 \BookCite{36}. 它们是离散化空间 $H^1$ 最自然的有限元. 空间 $\mathcal P_r\Lambda^n$ 与 $\mathcal P_{r+1}^-\Lambda^n$ (注意指标的变化) 也彼此重合, 并与所有次数至多为 $r$ 的分片多项式组成的更大空间等同, 无论这些多项式是否连续; 这是 $L^2$ 的自然有限元离散化. Reed 和 Hill 于 1973 年开创的间断 Galerkin 方法 \BookCite{65} 在有限元计算中使用这样的空间. 二维中的空间 $\mathcal P_r^-\Lambda^1$ 与 Raviart 和 Thomas 于 1975 年引入的著名空间 \BookCite{64} 等同, 用于离散化 $H(\operatorname{div})$. Nédélec 于 1980 年 \BookCite{59} 以两种不同方式将该空间推广至三维, 一种作为 $H(\operatorname{div})$ 的离散化, 另一种作为 $H(\operatorname{curl})$ 的离散化. 它们分别对应于三维中的空间 $\mathcal P_r^-\Lambda^1$ 和 $\mathcal P_r^-\Lambda^2$. 最后, 二维空间 $\mathcal P_r\Lambda^1$ 与 1985 年的 Brezzi–Douglas–Marini 空间等同, 而三维空间 $\mathcal P_r\Lambda^1$ 和 $\mathcal P_r\Lambda^2$ 则与 Nédélec 于 1986 年给出的推广等同, 称为第二类 Nédélec 棱单元和面单元. 因此, 在 $n\leq3$ 维时, 所有这些空间都已出现于有限元文献中. FEEC 的处理揭示了这些在不同背景下产生的空间之间的统一性. 在某些情形中, 这些空间最自然的自由度在 FEEC 发展以前还不为人所知.}

% block: geometric-and-feec-historical-antecedents page:104 type:ordinary-paragraph
{还值得注意的是, 这些分片多项式微分形式空间中, 有些独立地出现于几何背景下. 空间 $\mathcal P_1^-\Lambda^k$ 具有很强的几何性, 剖分中每个 $k$-单纯形对应一个自由度, 因而与单纯 $k$-上链空间同构. 事实上, Whitney 在 1957 年的著作 \BookCite{70} 中就构造了这些空间, 比它们作为混合有限元出现早数十年. 它们与混合有限元的联系则迟至 1988 年才由 Bossavit \BookCite{20} 指出. 此外, 与剖分关联的完全多项式空间 $\mathcal P_r\Lambda^k$ 出现在 Sullivan 于 1970 年代的拓扑学工作 \BookCite{67}, \BookCite{68} 中, 比它们作为 Brezzi–Douglas–Marini 单元被重新发现早十年. 在 1980 年代初的一本大体上被忽视的会议论文集 \BookCite{17} 中, Baker 将其称为 Sullivan–Whitney 形式, 并用它们离散化 Hodge Laplacian 的一个特征值问题, 分析了该离散化的收敛性. Hiptmair 于 1999 年的一篇奠基性论文 \BookCite{50} 首次统一处理了 $\mathcal P_r^-\Lambda^k$ 空间. 在 Arnold, Falk 和 Winther 于 2006 年发表的论文 \BookCite{11} 中, \emph{有限元外微积分}这一术语首次出现, Koszul 复形首次应用于有限元, 从而使 $\mathcal P_r^-\Lambda^k$ 和 $\mathcal P_r\Lambda^k$ 两类空间得以大幅简化并得到统一处理.}

% block: cubical-finite-element-antecedents page:104 type:ordinary-paragraph
{对于立方体有限元空间, 我们简要介绍了 $\mathcal Q_r^-\Lambda^k(\mathcal T_h)$ 和 $\mathcal S_r\Lambda^k(\mathcal T_h)$ 两族. 前者可视为 $\mathcal P_r^-\Lambda^k(\mathcal T_h)$ 空间的张量积对应物, 它们曾与相应的单纯形空间一同被导出. 例如, 二维中的 $\mathcal Q_r^-\Lambda^1$ 空间就是 \BookCite{64} 中的四边形 Raviart–Thomas 空间. $\mathcal S_r\Lambda^k$ 族则不同, 它在既有文献中基本未曾出现, 其发现受到 FEEC 的启发. 对 0-形式的单元, 即通常的 $H^1$ 有限元, 在二维中已广为人知, 就是常用的 Serendipity 有限元. $\mathcal S_r\Lambda^0$ 的构造将这些单元系统地推广到了更高维. 三维及更高维中用于 1-形式和 2-形式的 $\mathcal S_r$ 空间, 此前同样不为人所知.}

% block: finite-element-forms-summary-heading page:105 type:ordinary-paragraph

\refstepcounter{booksummary}\label{box:finite-element-forms-summary}
\begin{MBSourceBox}
\textbf{Box~\ref{box:finite-element-forms-summary}. {Chapter~\ref{chap:finite-element-differential-forms} 有限元微分形式总结.}}

% block: finite-element-forms-summary-spaces page:105 type:ordinary-paragraph
{给定任意 $n\geq1$ 的 $n$ 维单纯形剖分 $\mathcal T_h$, 对任意形式次数 $0\leq k\leq n$ 和任意多项式次数 $r\geq1$, 都有两个典范的有限元微分 $k$-形式空间, 由次数不超过 $r$ 的分片多项式组成, 即完全多项式空间 $\mathcal P_r\Lambda^k(\mathcal T_h)$ 和修剪多项式空间 $\mathcal P_r^-\Lambda^k(\mathcal T_h)$. 与任何有限元空间一样, 可以通过给出形函数和自由度并证明其唯一可解性来指定这些空间.}

% block: finite-element-forms-summary-trimming page:105 type:ordinary-paragraph
{顾名思义, 完全多项式空间的形函数由系数为次数至多 $r$ 的多项式的所有 $k$-形式组成. 修剪多项式空间将系数多项式限制在所有次数至多 $r$ 的多项式的某个子空间中, 但该子空间至少包含所有次数至多 $r-1$ 的多项式. 它借助 Koszul 复形精确定义. 对 0-形式, 修剪空间与相同次数的完全多项式空间重合; 对 $n$-形式, 则与次数为 $r-1$ 的完全多项式空间重合. 在所有其他情形中, 修剪空间不是完全多项式空间. 在最低阶情形 $r=1$ 时, 修剪空间就是 Whitney 形式空间.}

% block: finite-element-forms-summary-complexes page:105 type:ordinary-paragraph
{完全多项式空间在 $k$ 增加时让 $r$ 递减, 组合成具有有界上链投影的 $L^2$ de Rham 复形的子复形. 修剪多项式空间也类似地构成 de Rham 子复形, 但次数 $r$ 保持不变.}

% block: finite-element-forms-summary-methods page:105 type:ordinary-paragraph
{如 \eqref{eq:four-galerkin-form-pairs} 所示, 两个这样的空间, 一个用于 $k$-形式, 一个用于 $(k-1)$-形式, 无论两者均为完全空间, 均为修剪空间, 还是各取其一, 都可用于求解 $k$-形式 Hodge Laplacian.}

% block: finite-element-forms-summary-cubical page:105 type:ordinary-paragraph
{对于立方体网格而非单纯形剖分, 也存在类似的空间族.}
\end{MBSourceBox}
\par\endgroup
% Original PDF 106 is a blank production leaf with printer marks, retained in the source.
